% Dibangkitkan otomatis oleh examples/run_verification.py
% (jangan mengedit manual — Bagian 4.9.2 aturan 2)

\begin{table}[htbp]\centering
\caption{Reproduksi Tabel 4.1 (hasil solver, empat model utuh).}
\label{tab:repro41}
\begin{tabular}{lcccccc}\hline
Model & Batang & Rank/DOF & $u_{\max}$ (mm) & $U_{\max}$ & $W$ (kNmm) & Cek \\\hline
M1_parallel_bars_1d & 2 & 1/1 & 1.000000 & 0.500000 & 15.000000 & PASS \\
M2_planar_triangle & 3 & 3/3 & 1.154460 & 0.442218 & 16.956287 & PASS \\
M3_square_one_diagonal & 5 & 5/5 & 5.935311 & 0.530330 & 86.139610 & PASS \\
M4_square_two_diagonals & 6 & 5/5 & 3.582277 & 0.320083 & 51.989854 & PASS \\
\hline\end{tabular}\end{table}

\begin{table}[htbp]\centering
\caption{Reproduksi Tabel 4.3 (persegi satu diagonal, $\lambda=0$).}
\label{tab:repro43}
\begin{tabular}{lcccc}\hline
Batang & RR & CI & Rank/DOF & Cek \\\hline
$e_1$ & 0.983152 & 0.016848 & 4/5 & PASS \\
$e_2$ & 0.964724 & 0.216162 & 4/5 & PASS \\
$e_3$ & 1.000000 & 0.000000 & 4/5 & PASS \\
$e_4$ & 1.000000 & 0.000000 & 4/5 & PASS \\
$e_5$ & 0.717285 & 0.472914 & 4/5 & PASS \\
\hline\end{tabular}\end{table}

\begin{table}[htbp]\centering
\caption{Reproduksi Tabel 4.2 (status gangguan; kontrol tidak masuk jumlah gangguan).}
\label{tab:repro42}
\begin{tabular}{lcccc}\hline
Model & Gangguan & Solusi unik & Rank kurang & Terhubung \\\hline
M1_parallel_bars_1d & 8 & 8 & 0 & 8 (PASS) \\
M2_planar_triangle & 12 & 9 & 3 & 12 (PASS) \\
M3_square_one_diagonal & 20 & 15 & 5 & 20 (PASS) \\
M4_square_two_diagonals & 24 & 24 & 0 & 24 (PASS) \\
\hline\end{tabular}\end{table}

\begin{table}[htbp]\centering
\caption{Reproduksi Tabel 4.4 (dua kehilangan elemen bergaya awal nol).}
\label{tab:repro44}
\begin{tabular}{lccccc}\hline
Elemen & CI & $h_e$ & $m_e(0)$ & Rank & Gabungan \\
\hline
$e_3$ & 1.11e-16 & 1.000000 & 2.22e-16 & 4/5 & Teruskan \\
$e_4$ & 0.00e+00 & 1.000000 & 0.00e+00 & 4/5 & Teruskan \\
\hline\end{tabular}\end{table}

\begin{table}[htbp]\centering
\caption{Reproduksi Tabel 4.5 (paruh kanonik; ambang $\tau=0.5$, kontrol tidak dihitung).}
\label{tab:repro45}
\begin{tabular}{lcccc}\hline
Kebijakan & $N$ & Diteruskan & Rank terlewat & Kritis terlewat \\
\hline
ARI & 64 & 5 & 6/8 & 14/17 \\
Pengaman & 64 & 8 & 0/8 & 9/17 \\
Gabungan & 64 & 11 & 0/8 & 8/17 \\
\hline\end{tabular}\end{table}
